\chapter{Surfaces}\label{ch:surfaces}

Surfaces are the faces of every model: lens surfaces, the walls of a channel, the faces of a
solid. Like curves, all surfaces in \NF{} are instances of a single class,
\apiref{NurbsSurface}, a tensor-product NURBS. This chapter covers the exact constructions
(planar patches, extrusion, revolution, ruled surfaces), the free-form constructions through
curves (loft, sweeps, Coons and Gordon surfaces), the \code{NurbsSurface} class with all its
members, and the operations and analyses on surfaces (extension, offsets, fillets, curvature,
area). Surfaces carry no boundary of their own beyond their parameter rectangle; trimmed faces
are the subject of \cref{ch:trimming}, closed solids of \cref{ch:solids}.

\section{Definition and conventions}\label{sec:surface-def}

A NURBS surface of degrees $(p, q)$ with an $(n+1)\times(m+1)$ net of control points $\mathbf
P_{ij}$, weights $w_{ij} > 0$ and knot vectors $U$ (size $n+p+2$) and $V$ (size $m+q+2$) is
\[
  \mathbf S(u,v) = \frac{\sum_{i=0}^{n}\sum_{j=0}^{m} N_{i,p}(u)\,N_{j,q}(v)\,w_{ij}\,\mathbf P_{ij}}
                        {\sum_{i=0}^{n}\sum_{j=0}^{m} N_{i,p}(u)\,N_{j,q}(v)\,w_{ij}} .
\]
The index $i$ runs along $u$, $j$ along $v$. The net is stored \emph{row-major}: control point
$\mathbf P_{ij}$ has the flat index $i\cdot(m+1)+j$, i.e.\ $j$ (the $v$ index) varies fastest.

The orientation of a surface is given by its parametrisation: the (unnormalised) normal is
$\mathbf S_u \times \mathbf S_v$. Many operations depend on it -- offsets, fillets, the outward
direction of solid faces -- so it is worth knowing which way each construction points.
\Cref{tab:surface-params} summarises the parametrisations of the constructions in this chapter.

\begin{table}[htbp]
\centering\small
\begin{tabularx}{\linewidth}{@{}l>{\raggedright\arraybackslash}X>{\raggedright\arraybackslash}X@{}}
\toprule
Construction & $u$ & $v$ \\
\midrule
\code{make_rectangle(o, eu, ev)} & along \code{eu}, degree 1 & along \code{ev}, degree 1 \\
\code{make_bilinear(p00, p10, p01, p11)} & from \code{p00} to \code{p10} & from \code{p00} to \code{p01} \\
\code{extrude(profile, d)} & the profile & along \code{d}, degree 1 \\
\code{revolve(profile, ...)} & around the axis, rational degree 2 & the profile \\
\code{ruled(c0, c1)} & the curves (made compatible) & from \code{c0} ($v=0$) to \code{c1} ($v=1$), degree 1 \\
\code{loft(sections)} & along the sections & across the sections \\
\code{sweep}, \code{sweep2} & the profile, on $[0,1]$ & the trajectory (its domain) \\
\code{sweep_translational} & the profile & the trajectory \\
\code{coons(c0, c1, d0, d1)} & along \code{c0}, \code{c1} & along \code{d0}, \code{d1} \\
\bottomrule
\end{tabularx}
\caption{Parametrisation of the surface constructions.}
\label{tab:surface-params}
\end{table}

\section{Exact constructions}\label{sec:exact-surfaces}

\begin{description}[style=nextline]
\item[\apiref{make_rectangle}\code{(origin, edge_u, edge_v)}]
  The planar patch $\mathbf S(u,v) = \mathbf o + u\,\mathbf e_u + v\,\mathbf e_v$ on $[0,1]^2$,
  bilinear with $2\times2$ control points. Its normal is $\mathbf e_u \times \mathbf e_v$. The
  edges need not be perpendicular (any parallelogram). Rectangles are the usual carriers of
  planar faces: a disc is a square trimmed by a circle (\cref{ch:trimming}).
\item[\apiref{make_bilinear}\code{(p00, p10, p01, p11)}]
  The bilinear patch through four corners; with non-coplanar corners it is a hyperbolic
  paraboloid (a ``saddle''). \code{p10} is the corner at $(u,v)=(1,0)$, \code{p01} the one at
  $(0,1)$.
\item[\apiref{extrude}\code{(profile, direction)}]
  The general cylinder obtained by moving the profile along the vector \code{direction}.
  Exact for any profile, including rational ones; a circle extruded gives an exact cylinder.
\item[\apiref{revolve}\code{(profile, axis_point, axis_dir, angle)}]
  The surface of revolution (Piegl \& Tiller A8.1): the profile is rotated by \code{angle}
  (radians, $0 < \text{angle} \le 2\pi$) about the axis through \code{axis_point} with direction
  \code{axis_dir}, counter-clockwise by the right-hand rule. $u$ runs around the axis (rational
  quadratic, $\le 90^\circ$ per segment), $v$ follows the profile. A semicircle gives an exact
  sphere, a circle off the axis an exact torus, a line an exact cone or cylinder.
\item[\apiref{ruled}\code{(c0, c1)}]
  The ruled surface (linear loft) between two curves: straight lines connect points of equal
  parameter. The curves are made compatible automatically (\apiref{make_compatible}), so their
  parametrisations should correspond -- two circles should start at corresponding points and
  run in the same direction, otherwise the surface is twisted.
\end{description}

\begin{note}[Poles]
A profile point on the axis of revolution (the end points of the semicircle that makes a sphere)
produces a \emph{pole}: a whole row of the control net collapses into one point. There
$\mathbf S_u = 0$ and the usual normal $\mathbf S_u\times\mathbf S_v$ is undefined.
\code{normal()} then returns the \emph{limit normal}, computed from the mixed derivative, with
the sign chosen to agree with the neighbourhood. The example prints $(0,0,-1)$ at the south pole
of the sphere -- the outward normal, as on the rest of the sphere.
\end{note}

\example{06_basic_surfaces.py}

Output:

\begin{lstlisting}[style=text]
rectangle: degrees (1, 1), net 2 x 2, S(0.5, 0.5) = Vec3(1.000000, 0.500000, 0.000000), normal Vec3(0.000000, 0.000000, 1.000000), area 2.0000000000
bilinear: S(0.5, 0.5) = Vec3(0.500000, 0.500000, 0.250000), area 1.2807892753
extrude: degrees (2, 1), net 9 x 2, area 12.5663706144 (4 pi = 12.5663706144)
revolve: degrees (2, 2), net 9 x 5, area 12.5663706144 (4 pi)
  normal at the pole (limit normal): Vec3(-0.000000, 0.000000, -1.000000)
  normal at the equator: Vec3(1.000000, -0.000000, 0.000000)
torus area 39.47841760 (4 pi^2 R r = 39.47841760)
ruled: S(u, 1) radius 0.500000000000, area 5.26861105 (pi (r0 + r1) s = 5.26861105)
\end{lstlisting}

All areas agree with the analytic values to the printed digits -- the constructions are exact
and \apiref{area} integrates to a relative accuracy of $10^{-10}$.

\section{Free-form surfaces through curves}\label{sec:freeform-surfaces}

\subsection{Loft (skinning)}

\begin{lstlisting}[style=text]
loft(sections, degree_v=3, method=Parametrization.Centripetal) -> NurbsSurface
loft(sections, degree_v, v_params) -> NurbsSurface
loft_parameters(sections, method=Parametrization.Centripetal) -> list[float]
\end{lstlisting}

\apiref{loft} creates the surface through a sequence of section curves (Piegl \& Tiller,
section~10.3). The sections are first made compatible -- clamped, domain $[0,1]$, common
(highest) degree and the union of all knots -- and every column of control points is then
interpolated across the sections with degree $\min(\text{degree\_v}, K-1)$ for $K$ sections.
$u$ runs along the sections, $v$ across them, and $\mathbf S(u, v_k) = \mathbf C_k(u)$ exactly
for the section parameters $v_k$.

The section parameters are computed by \apiref{loft_parameters} from the distances between the
control polygons, averaged over all columns (\code{Centripetal} by default, \code{ChordLength}
or \code{Uniform} on request; see \cref{ch:fitting}). The first is 0 and the last 1. Alternatively
pass the parameters explicitly as a strictly increasing list \code{v_params}; the $v$ domain
then is $[v_0, v_{K-1}]$. Explicit parameters are useful when the sections belong to known
heights or times.

Rational sections (circles) are interpolated in homogeneous space. Linear data are reproduced
exactly: equal circles give a cylinder, linearly growing circles a cone. A loft needs at least
two sections; a non-positive weight in the interpolated net raises \code{GeometryError}.

\begin{warning}[Corresponding sections]
The sections must be parametrised alike: same start point (seam) and same direction. Two circles
created with \code{make_circle} and parallel normals do correspond. Mixing a circle with a
polygon, or circles with opposite normals, produces a twisted surface. Use
\code{make_arc} with explicit axes or \apiref{reparametrize}/\code{reversed} to align sections.
\end{warning}

\subsection{Sweeps}

\begin{lstlisting}[style=text]
sweep(profile, trajectory, options=SweepOptions()) -> NurbsSurface
sweep2(profile, rail1, rail2, options=SweepOptions()) -> NurbsSurface
sweep_translational(profile, trajectory) -> NurbsSurface
SweepOptions(frame=SweepFrame.RotationMinimizing, twist=0.0, tolerance=1e-6, max_sections=1025)
\end{lstlisting}

\apiref{sweep} moves a profile along a trajectory. The profile is given \emph{in world
coordinates at the start of the trajectory}; at the trajectory parameter $v$ it is translated to
$\mathbf T(v)$ and rotated by the change of the moving frame from $\mathbf T(v_0)$ to
$\mathbf T(v)$. The \apiref{SweepOptions} control the result:

\begin{tabularx}{\linewidth}{@{}l l >{\raggedright\arraybackslash}X@{}}
\toprule
Field & Default & Meaning \\
\midrule
\code{frame} & \code{RotationMinimizing} & \code{SweepFrame.RotationMinimizing}: rotation-minimising
  frame (double reflection method of Wang, Jüttler, Zheng and Liu); no spurious twist, defined
  on straight parts and at inflections. \code{SweepFrame.Fixed}: pure translation without rotation. \\
\code{twist} & 0 & additional rotation about the tangent at the end of the trajectory (radians),
  growing linearly with the trajectory parameter \\
\code{tolerance} & $10^{-6}$ & allowed deviation from the ideal swept surface (model units) \\
\code{max_sections} & 1025 & upper limit of placed profile copies \\
\bottomrule
\end{tabularx}

The swept surface is not exactly a NURBS in general. \NF{} places exact copies of the profile
along the trajectory and skins them with \code{loft}, piecewise between the knots of the
trajectory; the number of copies is doubled until the deviation from the ideal sweep, measured
between the copies, is below the tolerance. \code{GeometryError} is raised if
\code{max_sections} is not enough or the trajectory has a vanishing tangent. With
\code{SweepFrame.Fixed} and no twist the exact translational sweep is returned.

\apiref{sweep2} is the two-rail sweep: the profile, given at the starts of both rails, is moved,
rotated and \emph{uniformly scaled} so that its points on \code{rail1} and \code{rail2} follow
the rails. \code{rail2} is linearly reparametrised to the domain of \code{rail1}; the rotation
about the rail-to-rail axis follows the tangent of \code{rail1}. In the example the arch between
two diverging lines doubles its span and its height.

\apiref{sweep_translational} is exact: $\mathbf S(u,v) = \mathbf C(u) + \mathbf T(v) -
\mathbf T(v_0)$, a tensor product of the two curves (rational curves are allowed).

\paragraph{Several profiles.} \code{sweep(profiles, trajectory, params=[], options)} takes a
list of profiles, each drawn \emph{where it crosses the trajectory}, at the trajectory
parameters \code{params} (empty: the parameter of the trajectory point nearest to each
profile). Every profile is moved back to the trajectory start by the inverse of its placement,
the moved profiles are lofted over their parameters (cubic for four or more profiles), and the
section at $v$ is that blended shape placed at $v$. Before the first and after the last profile
the shape is constant. Closed profiles should have their seams aligned, as for \code{loft}.

\paragraph{Rail revolve.} \apiref{rail_revolve} revolves a profile about an axis while the
profile end nearer to a \emph{rail} curve follows the rail:

\begin{lstlisting}[style=text]
rail_revolve(profile, rail, axis_point, axis_dir, scale_height=False, options=SweepOptions()) -> NurbsSurface
\end{lstlisting}

At the rail parameter $v$ the profile is rotated about the axis to the angle of
$\mathbf R(v)$ and scaled about the axis in the radial direction, so that its end lies at the
rail's distance from the axis. With \code{scale_height} it is also scaled along the axis (about
the axis point) to the rail's height; otherwise heights are kept. A profile starting on the axis
keeps that point, e.g.\ a reflector with an elliptic or rounded-square rim. $u$ runs along the
profile, $v$ along the rail; a closed rail gives a closed surface.


\subsection{Coons patches and Gordon surfaces}

\begin{lstlisting}[style=text]
coons(c0, c1, d0, d1, tolerance=1e-7) -> NurbsSurface
gordon(u_curves, v_curves, tolerance=1e-7) -> NurbsSurface
\end{lstlisting}

A \apiref{coons} patch fills a closed loop of four boundary curves: $\mathbf S(u,0) =
\mathbf c_0$, $\mathbf S(u,1) = \mathbf c_1$, $\mathbf S(0,v) = \mathbf d_0$, $\mathbf S(1,v) =
\mathbf d_1$. The curves must meet at the four corners; their directions are fixed
automatically. The interior is the bilinearly blended Boolean sum of the two ruled surfaces minus
the bilinear corner patch.

A \apiref{gordon} surface interpolates a whole \emph{network} of curves: a family of
$u$-curves and a family of $v$-curves that intersect each other. The network must be
\emph{consistent}: every $u$-curve must intersect every $v$-curve, and all intersection points of
one $u$-curve must lie at the same $v$-curve parameters as on the others. The intersection
parameters are determined and checked internally. The Gordon surface is again a Boolean sum
$L_1 + L_2 - T$ of two lofts and a tensor-product interpolant of the intersection points.

The Boolean sum is linear in the control points. For non-rational curves it is therefore
exact. For rational boundaries (arcs, conics) \code{coons} forms it in homogeneous coordinates
-- numerators $w\mathbf P$ and weights $w$ separately -- which reproduces the boundaries
exactly if the weights agree at the corners. Each curve's weights may be scaled for this
(scaling does not change a curve); arcs from \code{make_arc}, whose end weights are 1, always
qualify. A Coons patch from two coaxial arcs and two straight lines is then an exact cylinder
patch. Otherwise, and always for \code{gordon}, rational input is first replaced by cubic
interpolants within \code{tolerance}.

\apiref{triangle_patch} fills a loop of three curves: a Coons patch with \code{c0} as one side
and the opposite corner as a degenerate side ($v = 1$), so it is regular except at that corner,
where the limit normal is used:

\begin{lstlisting}[style=text]
triangle_patch(c0, c1, c2, tolerance=1e-7) -> NurbsSurface
\end{lstlisting}

\subsection{Example}

\example{06_freeform_surfaces.py}

Output:

\begin{lstlisting}[style=text]
loft: net 9 x 4, degree_v 3, section parameters [0.0, 0.3172, 0.6987, 1.0]
  surface passes through the sections: max radius error 2.2e-16
  explicit v_params: domain_v [0.0, 3.0], S(0, 2) = Vec3(0.000000, 0.800000, 2.500000)
sweep: net 9 x 65, max deviation from tube radius 1.4e-08
  area 24.674011 (torus section 2 pi r * R pi/2 = 24.674011)
twisted sweep: corner at the top Vec3(0.500000, -0.500000, 3.000000)
translational sweep: degrees (3, 2), S(0, 1) - S(0, 0) = Vec3(0.000000, -0.244835, 0.958851)
sweep2: span at the start 1.000000, at the end 2.000000, top at the end 1.000000
coons: S(0.5, 0) = Vec3(1.000000, 0.000000, 1.000000), S(0.5, 0.5) = Vec3(1.000000, 1.000000, 0.500000), S(0.5, 1) = Vec3(1.000000, 2.000000, 0.000000)
gordon: net 9 x 3, network -> surface distance 3.2e-13
\end{lstlisting}

The bent tube lies within $1.4\cdot10^{-8}$ of radius 0.5 around its path, and its area agrees
with the analytic value of the quarter torus. The twisted square's first corner, which starts at
$(-0.5,-0.5,0)$, ends rotated by $90^\circ$ at $(0.5,-0.5,3)$.

\section{The NurbsSurface class}\label{sec:nurbssurface}

\subsection{Construction and properties}

\begin{lstlisting}[style=text]
NurbsSurface(degree_u, degree_v, knots_u, knots_v, num_u, num_v, control_points, weights=[])
\end{lstlisting}

\code{control_points} is an array of shape $(\text{num\_u}\cdot\text{num\_v}, 3)$ in row-major
order (index $i\cdot\text{num\_v}+j$, $i$ along $u$); \code{weights} has the same order. The
constructor validates sizes, knot vectors and weights.

\begin{tabularx}{\linewidth}{@{}l>{\raggedright\arraybackslash}X@{}}
\toprule
Property & Meaning \\
\midrule
\code{degree_u}, \code{degree_v} & degrees $p$ and $q$ \\
\code{knots_u}, \code{knots_v} & knot vectors \\
\code{num_u}, \code{num_v} & size of the control net \\
\code{control_points} & array of shape $(\text{num\_u}\cdot\text{num\_v}, 3)$ \\
\code{weights} & list of weights in the same order \\
\code{domain_u}, \code{domain_v} & parameter ranges as two-element lists \\
\bottomrule
\end{tabularx}

\subsection{Evaluation}

\begin{description}[style=nextline]
\item[\code{s.evaluate(u, v)} or \code{s(u, v)}] the point $\mathbf S(u,v)$.
\item[\code{s.first_derivatives(u, v)}] the list $[\mathbf S, \mathbf S_u, \mathbf S_v]$ (fast path).
\item[\code{s.derivatives(u, v, order=1)}] the nested list \code{D} with
  \code{D[k][l]} $= \partial^{k+l}\mathbf S/\partial u^k\partial v^l$ for $k+l \le
  \text{order}$ (Piegl \& Tiller A3.6 and A4.4, any order, rational quotient rule included).
\item[\code{s.normal(u, v)}] the unit normal $\widehat{\mathbf S_u\times\mathbf S_v}$, the
  limit normal at poles; raises \code{GeometryError} if both partial derivatives vanish.
\item[\code{s.iso_curve_u(u)}, \code{s.iso_curve_v(v)}] the iso-parametric curves as
  \code{NurbsCurve}s: \code{iso_curve_u(u)} is the curve $v \mapsto \mathbf S(u,v)$ at fixed
  $u$, \code{iso_curve_v(v)} the curve $u \mapsto \mathbf S(u,v)$. They are exact.
\item[\code{s.sample(nu=20, nv=20)}] a NumPy array of shape $(\text{nu}, \text{nv}, 3)$ at equally
  spaced parameters including the boundaries -- ready for \code{plot_surface} in Matplotlib.
\end{description}

\subsection{Shape-preserving operations}

Every curve operation of \cref{sec:nurbscurve} has a surface counterpart per direction:

\begin{tabularx}{\linewidth}{@{}l>{\raggedright\arraybackslash}X@{}}
\toprule
Method & Effect \\
\midrule
\code{refine_knots_u(knots)}, \code{refine_knots_v} & insert knots (A5.5) \\
\code{try_remove_knot_u(u, times=1, tolerance=1e-10)}, \code{..._v} & remove a knot as often as
  possible within tolerance; returns \code{(surface, removed)}. The same number of removals is
  applied to all rows. \\
\code{remove_knot_u(u, times=1, tolerance=1e-10)}, \code{..._v} & as above, raising
  \code{GeometryError} if not possible \\
\code{elevate_degree_u(times=1)}, \code{..._v} & degree elevation (A5.10) \\
\code{reduce_degree_u(tolerance=1e-10)}, \code{..._v} & degree reduction by one, verified \\
\code{split_u(u)}, \code{split_v(v)} & two surfaces keeping their parameter intervals \\
\code{sub_surface(u0, u1, v0, v1)} & the piece over a parameter rectangle \\
\code{clamped()} & clamped knot vectors in both directions \\
\code{decompose_bezier()} & Bézier patches, row-major ($u$ span outer), on their original
  parameter rectangles \\
\code{transposed()} & swaps $u$ and $v$: $\mathbf T(v,u) = \mathbf S(u,v)$; this reverses the
  normal \\
\code{transformed(transform)} & affine image (\cref{ch:transforms}) \\
\bottomrule
\end{tabularx}

\apiref{join_u}\code{(a, b, tolerance=1e-10)} appends surface \code{b} to \code{a} across
\code{a}'s $u = u_{\max}$ boundary, \apiref{join_v} across $v = v_{\max}$. The boundary rows
must agree \emph{as control rows}, i.e.\ both surfaces must share the parametrisation of the
common boundary and not only its geometry; the $v$ degrees and knots are made compatible.
\code{join_u(*s.split_u(u))} reproduces the surface.

\subsection{Example}

\example{06_nurbs_surface.py}

Output:

\begin{lstlisting}[style=text]
degrees (3, 3), net 4 x 4, domain [0.0, 1.0] x [0.0, 1.0]
S(0.5, 0.5) = Vec3(1.500000, 1.000000, 0.562500)
S_u = Vec3(3.000000, 0.000000, 1.125000)  S_v = Vec3(0.000000, 2.000000, 0.000000)
S_uv = Vec3(0.000000, 0.000000, 0.000000)  S_vv = Vec3(0.000000, 0.000000, -3.375000)
unit normal Vec3(-0.351123, 0.000000, 0.936329)
sample shape (5, 3, 3) corner [3. 2. 0.]
iso_curve_u(0.5)(0.5) = Vec3(1.500000, 1.000000, 0.562500), iso_curve_v(0.5)(0.5) = Vec3(1.500000, 1.000000, 0.562500)
refined/elevated: net 7 x 5, degree_v 4, deviation 9.2e-16
knot u = 0.5 removed 1 time(s); degree_v back to 3
split_u domains [0.0, 0.4] [0.4, 1.0]
sub_surface domain [0.1, 0.9]
Bezier patches of the refined surface: 4
transposed: T(v, u) == S(u, v): True
join_u of the halves: deviation 4.4e-16
sphere r = 2: K = 0.2500000000, H = -0.5000000000, k1 = -0.4999999947, k2 = -0.5000000053
area of the upper half: 25.1327412287 (8 pi = 25.1327412287)
closest point: (u, v) = (0.000000, 0.792893), point Vec3(1.200000, 0.000000, 1.600000), distance 3.000000
integral of sin over [0, pi] = 2.0000000000000004
integral of sqrt(|x|) over [-1, 1] with a break at 0 = 1.333333333337871
\end{lstlisting}

\section{Analysis}\label{sec:surface-analysis}

\begin{description}[style=nextline]
\item[\apiref{surface_curvature}\code{(surface, u, v)}] returns a \apiref{SurfaceCurvature} with
  the Gaussian curvature \code{gaussian} $K$, the mean curvature \code{mean} $H$, the principal
  curvatures \code{k1}, \code{k2}, the principal directions \code{dir1}, \code{dir2} and the unit
  \code{normal}. The values come from the first and second fundamental forms with $\mathbf n =
  \widehat{\mathbf S_u\times\mathbf S_v}$; the signs of $H$, $k_1$, $k_2$ follow this
  orientation (negative for a sphere with outward normal: the surface bends away from the
  normal). At poles the values are extrapolated from nearby regular points. At umbilic points
  (sphere, plane) the principal directions are arbitrary and $k_1$, $k_2$ are slightly noisy,
  as the example shows.
\item[\apiref{area}\code{(surface, rel_tol=1e-10)}] the area of the whole parameter rectangle;
  \code{area(surface, u0, u1, v0, v1, rel_tol)} that of a sub-rectangle. Adaptive tensor
  Gauss--Kronrod quadrature per knot rectangle. For trimmed faces use
  \code{TrimmedSurface.area()} (\cref{ch:trimming}).
\item[\apiref{closest_point}\code{(surface, point)}] and
  \apiref{closest_point_local}\code{(surface, point, u0, v0)} as for curves; they return a
  \apiref{SurfacePoint} with \code{u}, \code{v}, \code{point} and \code{distance}.
\item[\apiref{integrate}\code{(f, a, b, rel_tol=1e-10, breaks=[], abs_tol=0.0)}] the globally
  adaptive Gauss--Kronrod quadrature used internally, available for any Python function of one
  variable. \code{breaks} lists interior points where the integrand is not smooth (kinks,
  jumps); the interval is split there.
\end{description}

\section{Extending surfaces}\label{sec:surface-extend}

\begin{lstlisting}[style=text]
extend(surface, side, length, type=Extension.Natural) -> NurbsSurface
\end{lstlisting}

\apiref{extend} extends a surface across one of its four sides, \code{SurfaceSide.UMin},
\code{UMax}, \code{VMin} or \code{VMax}, by an arc length measured along the middle iso-curve.
The original parameters are kept. \code{Extension.Natural} continues the surface polynomially
(quadrics stay quadrics -- a cylinder band extended around its axis stays on the cylinder);
\code{Extension.Linear} continues along the tangent planes, which requires constant weights along
the iso-curves across the side.

\section{Offset surfaces}\label{sec:offset-surface}

\begin{lstlisting}[style=text]
offset_surface(surface, distance, tolerance=1e-7) -> NurbsSurface
\end{lstlisting}

\apiref{offset_surface} returns $\mathbf O(u,v) = \mathbf S(u,v) + d\,\mathbf N(u,v)$ with the
unit normal $\mathbf N$ (the limit normal at poles), i.e.\ positive distances go along
$\mathbf S_u\times\mathbf S_v$. Check the direction with \code{surface.normal(u, v)} before
offsetting. The result is a non-rational degree-5 tensor-product spline fitted to the tolerance;
the parametrisation is kept. \code{GeometryError} is raised where $d$ reaches a principal radius
of curvature (the offset would degenerate) and at creases (normal discontinuities across knot
lines) that would open a gap. Global self-intersections are not detected.

\section{Rolling-ball fillets}\label{sec:surface-fillet}

\begin{lstlisting}[style=text]
fillet_surfaces(s1, s2, radius, side1=1, side2=1, tolerance=1e-6,
                offset_tolerance=0.0, max_sections=1025) -> SurfaceFillet
\end{lstlisting}

\apiref{fillet_surfaces} computes the constant-radius fillet that a ball of the given radius
traces when it rolls in contact with both surfaces. \code{side1 = +1} places the ball on the side
of the normal $\mathbf S_u\times\mathbf S_v$ of \code{s1}, \code{-1} on the other side; likewise
\code{side2}. Spine points are solved exactly; the sections are exact circular arcs skinned
cubically, and intervals are bisected until the fillet, the spine and the contact curves are
within \code{tolerance} of the exact sections. \code{offset_tolerance} controls the fit of the
offset surfaces used for the start values (0 = automatic) and \code{max_sections} limits the
number of sections.

The result \apiref{SurfaceFillet} contains:

\begin{tabularx}{\linewidth}{@{}l>{\raggedright\arraybackslash}X@{}}
\toprule
Field & Meaning \\
\midrule
\code{surface} & the fillet surface; $u$ runs across the fillet from contact 1 to contact 2, $v$
  along the spine \\
\code{spine} & the curve of ball centres \\
\code{contact1}, \code{contact2} & the contact curves on \code{s1} and \code{s2} (3D) \\
\code{pcurve1}, \code{pcurve2} & the contact curves in the parameter spaces of \code{s1} and
  \code{s2}, as curves $(u_i(t), v_i(t), 0)$ -- ready to trim the input surfaces (\cref{ch:trimming}) \\
\bottomrule
\end{tabularx}

Not handled: several spine branches (the longest intersection curve of the offsets is used).
The input surfaces are not trimmed automatically; \apiref{fillet_edges} does this for the edges
of a solid, including vertex blends where several fillets meet and variable radii
(\cref{sec:edge-blends}). In C++, \code{SurfaceFilletOptions::radius_at} gives a variable
radius as a function of the ball position.

\example{06_surface_operations.py}

Output:

\begin{lstlisting}[style=text]
extend UMax: domain_u [0.0, 1.4038182881544163], end point Vec3(-0.479426, 0.877583, 0.500000), radius 1.000000000000
extend VMax: domain_v [0.0, 1.5], top z = 1.500000
sphere normal at the equator Vec3(0.813826, 0.581109, -0.000000) (points outward)
offset +0.5: degrees (5, 5), net 49 x 25, radius 1.5 within 1.3e-07
offset -1.2 refused: offset_surface: |distance| = 1.200000 reaches the radius of curvature  ...
fillet: net 3 x 9, spine at mid Vec3(0.500000, 1.500000, 0.500000)
  contact on floor Vec3(0.500000, 1.500000, 0.000000), on wall Vec3(-0.000000, 1.500000, 0.500000)
  p-curve on the floor (u, v) = (0.300000, 0.500000) -> Vec3(0.500000, 1.500000, 0.000000)
  fillet surface is the quarter cylinder r = 0.5 within 3.3e-16
\end{lstlisting}

The fillet between the floor $z=0$ and the wall $x=0$ is the exact quarter cylinder of radius
0.5 around the line $x = z = 0.5$; the p-curve on the floor gives the parameters $(0.3, 0.5)$,
which map back to the contact point.

\section{Construction helpers}\label{sec:surface-tools}

\begin{lstlisting}[style=text]
extrude_tapered(profile, direction, angle, tolerance=1e-7) -> list[NurbsSurface]
ribbon(curve, width, plane_normal=(0, 0, 1), tolerance=1e-7) -> list[NurbsSurface]
extrude_to_point(curve, apex) -> NurbsSurface
rebuild_curve(curve, num_ctrl, degree=3) -> RebuiltCurve        # .curve, .max_error
rebuild_surface(surface, num_u, num_v, degree_u=3, degree_v=3) -> RebuiltSurface
\end{lstlisting}

\apiref{extrude_tapered} extrudes a planar profile along a direction normal to its plane, the
vector length being the height, with a draft angle: the top curve is the profile offset in its
plane by $h\tan\alpha$ (\code{offset_curve}), outwards for $\alpha > 0$, inwards for
$\alpha < 0$, and moved by \code{direction}. For an open profile ``outwards'' is its left side
seen from the tip of the direction. The profile is split at its corners; the offset pieces are
\emph{mitred} (extended or cut where neighbours meet), and each piece gives one ruled surface
from the profile piece to its top piece. Line pieces give exact planar faces, curved pieces are
within the offset tolerance. Together with \code{cap} this gives drafted solids, e.g.\ the
frustum in the example.

\apiref{ribbon} is the planar counterpart: ruled surfaces between a curve and its offset by
\code{width} in the plane with normal \code{plane_normal} (to the left for positive width),
mitred at the corners.

\apiref{extrude_to_point} is the cone over a curve: the ruled surface from the curve ($v = 0$)
to the apex ($v = 1$). The apex is a degenerate curve with the curve's degree, knots and
weights, so the rulings are straight also for rational curves and a circle gives an exact
circular cone.

\apiref{rebuild_curve} and \apiref{rebuild_surface} replace a curve or surface by a
non-rational B-spline with a prescribed number of control points, degrees and uniform knots on
the same domain, e.g.\ to unify surfaces before meshing or to simplify fitted data. The fit is
a least-squares fit of dense samples \emph{at their parameters}, so the parametrisation is kept;
\code{max_error} is the largest distance between input and result at equal parameters on a
grid four times denser than the fit samples, an upper bound of the geometric deviation. The
rational circle converges only slowly because its parametrisation is only $C^1$ at its knots.

\example{06_surface_tools.py}

Output:

\begin{lstlisting}[style=text]
duct: radius at z = 3: 1.500000, at the ends 1.000000
rail revolve: area 6.283185 (ellipse pi a b = 6.283185)
tapered extrusion: 4 faces, volume 4.746762862 (h/3 (A1 + A2 + sqrt(A1 A2)) = 4.746762862)
ribbon: 2 pieces, area 1.250000000 (1.25)
cone: area 7.024814731 (pi r s = 7.024814731)
coons from arcs: rational=True, max |r - 1| = 1.1e-16
triangle patch: boundary on the sphere within 1.1e-16, corner Vec3(0.000000, 0.000000, 1.000000)
rebuild circle with  8 control points: max error 1.82e-02
rebuild circle with 16 control points: max error 3.29e-03
rebuild circle with 32 control points: max error 5.82e-04
\end{lstlisting}
